% --- Archon physics formalization source begin ---
% archon:physics
% archon:covers PhyXMiniProblems/problem_phyx_mini_0044.lean
% archon:source-report reports/phyx_mini/problem_phyx_mini_0044.source.json

\chapter{Physics problem phyx\_mini\_0044}
\label{ch:PhyXMiniProblems_problem_phyx_mini_0044}

\paragraph{Problem source.}
\#\# Physical scenario

The near point of a certain hyperopic eye is 100 cm in front of the eye.

\#\# Question

Find the focal length of the contact lens that will permit the wearer to see clearly an object that is 25 cm in front of the eye.

\#\# Answer choices

- A: A: 35cm

- B: B: 33cm

- C: C: 37cm

- D: D: 43cm

\#\# Figure caption (auxiliary; use the image as primary evidence)

The image is a diagram illustrating an optical setup involving a converging lens. Here's a detailed breakdown of the content:

1. **Objects and Labels**:
   - **Converging Lens**: Positioned at the center, depicted by a double convex lens.
   - **Object**: Represented by a point labeled as "Object" with a small vertical arrow. It is located to the left of the lens.
   - **Image**: Shown as a vertical arrow on the left, labeled "Image."

2. **Distances and Measurements**:
   - The object distance from the lens is marked as \textbackslash{}( s = 25 \textbackslash{}, \textbackslash{}text\{cm\} \textbackslash{}).
   - The image distance is marked as \textbackslash{}( s' = -100 \textbackslash{}, \textbackslash{}text\{cm\} \textbackslash{}).
   - The focal length is labeled as \textbackslash{}( f \textbackslash{}) with no numerical value provided.

3. **Arrows and Lines**:
   - A straight horizontal line represents the optical axis.
   - Light rays travel from the object through the lens, converging towards the image. These rays are shown as solid and dashed lines.
   - Perpendicular lines are used to denote the measurements \textbackslash{}( s \textbackslash{}), \textbackslash{}( s' \textbackslash{}), and \textbackslash{}( f \textbackslash{}).

4. **Textual Elements**:
   - Labels include "Image," "Object," and "Converging lens."
   - Measurements are provided in centimeters, showing the relative positions and properties of the optical components.

The diagram illustrates the relationship between an object, the lens, and the resulting image, demonstrating concepts like magnification and focal distance in optics.

\#\# Dataset metadata

Geometrical Optics; Physical Model Grounding Reasoning, Multi-Formula Reasoning

\paragraph{Recorded answer/context.}
B: B: 33cm

\paragraph{Figure/image path.}
/root/proposal\_for\_physic/science-mango/phyx\_mini\_run/phyx\_data/test\_image/44.png

\paragraph{Formalization target.}
create a compiling Lean file with sorry bodies at `PhyXMiniProblems/problem\_phyx\_mini\_0044.lean`. The Lean declarations must preserve the physical quantities, dimensions or dimensional roles, figure labels, governing-law hypotheses, and final relation expressed by this problem.
Use Mathlib/Physlib names found through LeanExplore where available. If a physics API is missing, introduce faithful local abstractions rather than scalar placeholder aliases.

\begin{theorem}[Physics formalization target]
\label{thm:physics:phyx_mini_0044:target}
\lean{PhyXMiniProblems.ProblemPhyXMini0044.problem_phyx_mini_0044}
\uses{def:physics:phyx-mini-0044:phyxminiproblems-problemphyxmini0044-lengthincentimeters, def:physics:phyx-mini-0044:phyxminiproblems-problemphyxmini0044-contactlenscorrectionsetup, def:physics:phyx-mini-0044:phyxminiproblems-problemphyxmini0044-matchesscenarioreadouts, def:physics:phyx-mini-0044:phyxminiproblems-problemphyxmini0044-matchesprimaryfigure, def:physics:phyx-mini-0044:phyxminiproblems-problemphyxmini0044-usescontactlenseyeplanemodel, def:physics:phyx-mini-0044:phyxminiproblems-problemphyxmini0044-hasdepictedsignconfiguration, def:physics:phyx-mini-0044:phyxminiproblems-problemphyxmini0044-satisfiessignedthinlensequation, def:physics:phyx-mini-0044:phyxminiproblems-problemphyxmini0044-matchesanswertonearestcentimeter}
The assigned autoformalize agent should translate this physics problem into Lean declarations in the covered file, with theorem and lemma proof bodies written as `by sorry`.
\end{theorem}
\begin{proof}
This is an autoformalization task, not a proof task. Produce statements that can later be proved by the physics prover without weakening the physical model.
\end{proof}

% --- Archon declaration topology begin ---
\section{Declaration topology}
\begin{definition}[Length Quantity]
  \label{def:physics:phyx-mini-0044:phyxminiproblems-problemphyxmini0044-lengthquantity}
  \lean{PhyXMiniProblems.ProblemPhyXMini0044.LengthQuantity}
  A signed physical length whose numerical value changes coherently with units
\end{definition}
\begin{proof}
  This is the declared model definition of Length Quantity.
\end{proof}
\begin{definition}[centimeter Unit Choices]
  \label{def:physics:phyx-mini-0044:phyxminiproblems-problemphyxmini0044-centimeterunitchoices}
  \lean{PhyXMiniProblems.ProblemPhyXMini0044.centimeterUnitChoices}
  SI base units with centimeters selected as the length unit
\end{definition}
\begin{proof}
  This is the declared model definition of centimeter Unit Choices.
\end{proof}
\begin{definition}[length In Centimeters]
  \label{def:physics:phyx-mini-0044:phyxminiproblems-problemphyxmini0044-lengthincentimeters}
  \lean{PhyXMiniProblems.ProblemPhyXMini0044.lengthInCentimeters}
  \uses{def:physics:phyx-mini-0044:phyxminiproblems-problemphyxmini0044-lengthquantity, def:physics:phyx-mini-0044:phyxminiproblems-problemphyxmini0044-centimeterunitchoices}
  The signed scalar readout of a physical length in centimeters
\end{definition}
\begin{proof}
  This is the declared model definition of length In Centimeters.
\end{proof}
\begin{definition}[Thin Lens Kind]
  \label{def:physics:phyx-mini-0044:phyxminiproblems-problemphyxmini0044-thinlenskind}
  \lean{PhyXMiniProblems.ProblemPhyXMini0044.ThinLensKind}
  The two paraxial thin-lens types, distinguished physically by focal-length sign
\end{definition}
\begin{proof}
  This is the declared model definition of Thin Lens Kind.
\end{proof}
\begin{definition}[Geometrical Image Kind]
  \label{def:physics:phyx-mini-0044:phyxminiproblems-problemphyxmini0044-geometricalimagekind}
  \lean{PhyXMiniProblems.ProblemPhyXMini0044.GeometricalImageKind}
  Whether the geometrical image is formed by actual rays or by their backward extensions
\end{definition}
\begin{proof}
  This is the declared model definition of Geometrical Image Kind.
\end{proof}
\begin{definition}[Hyperopic Eye]
  \label{def:physics:phyx-mini-0044:phyxminiproblems-problemphyxmini0044-hyperopiceye}
  \lean{PhyXMiniProblems.ProblemPhyXMini0044.HyperopicEye}
  \uses{def:physics:phyx-mini-0044:phyxminiproblems-problemphyxmini0044-lengthquantity}
  The unaided eye datum relevant to the correction problem
\end{definition}
\begin{proof}
  This is the declared model definition of Hyperopic Eye.
\end{proof}
\begin{definition}[Thin Contact Lens]
  \label{def:physics:phyx-mini-0044:phyxminiproblems-problemphyxmini0044-thincontactlens}
  \lean{PhyXMiniProblems.ProblemPhyXMini0044.ThinContactLens}
  \uses{def:physics:phyx-mini-0044:phyxminiproblems-problemphyxmini0044-lengthquantity, def:physics:phyx-mini-0044:phyxminiproblems-problemphyxmini0044-thinlenskind}
  A thin contact lens and its signed focal length
\end{definition}
\begin{proof}
  This is the declared model definition of Thin Contact Lens.
\end{proof}
\begin{definition}[Contact Lens Correction Setup]
  \label{def:physics:phyx-mini-0044:phyxminiproblems-problemphyxmini0044-contactlenscorrectionsetup}
  \lean{PhyXMiniProblems.ProblemPhyXMini0044.ContactLensCorrectionSetup}
  \uses{def:physics:phyx-mini-0044:phyxminiproblems-problemphyxmini0044-lengthquantity, def:physics:phyx-mini-0044:phyxminiproblems-problemphyxmini0044-geometricalimagekind, def:physics:phyx-mini-0044:phyxminiproblems-problemphyxmini0044-hyperopiceye, def:physics:phyx-mini-0044:phyxminiproblems-problemphyxmini0044-thincontactlens}
  The physical quantities and figure labels in the correction setup. desiredObjectDistanceFromEye and the eye's near-point distance are positive magnitudes measured in front of the eye. objectDistanceFromLens and signedImageDistanceFromLens use the signed convention of the thin-lens equation. The contact-lens approximation places the lens at the eye plane
\end{definition}
\begin{proof}
  This is the declared model definition of Contact Lens Correction Setup.
\end{proof}
\begin{definition}[Matches Scenario Readouts]
  \label{def:physics:phyx-mini-0044:phyxminiproblems-problemphyxmini0044-matchesscenarioreadouts}
  \lean{PhyXMiniProblems.ProblemPhyXMini0044.MatchesScenarioReadouts}
  \uses{def:physics:phyx-mini-0044:phyxminiproblems-problemphyxmini0044-lengthincentimeters, def:physics:phyx-mini-0044:phyxminiproblems-problemphyxmini0044-contactlenscorrectionsetup}
  Numerical distances stated in the problem text, before the eye-plane approximation
\end{definition}
\begin{proof}
  This is the declared model definition of Matches Scenario Readouts.
\end{proof}
\begin{definition}[Matches Primary Figure]
  \label{def:physics:phyx-mini-0044:phyxminiproblems-problemphyxmini0044-matchesprimaryfigure}
  \lean{PhyXMiniProblems.ProblemPhyXMini0044.MatchesPrimaryFigure}
  \uses{def:physics:phyx-mini-0044:phyxminiproblems-problemphyxmini0044-lengthincentimeters, def:physics:phyx-mini-0044:phyxminiproblems-problemphyxmini0044-contactlenscorrectionsetup}
  Readouts and qualitative labels transcribed from the primary figure: the contact lens is converging, s = 25 cm, and the backward ray extensions form a virtual image at signed distance s' = -100 cm
\end{definition}
\begin{proof}
  This is the declared model definition of Matches Primary Figure.
\end{proof}
\begin{definition}[Uses Contact Lens Eye Plane Model]
  \label{def:physics:phyx-mini-0044:phyxminiproblems-problemphyxmini0044-usescontactlenseyeplanemodel}
  \lean{PhyXMiniProblems.ProblemPhyXMini0044.UsesContactLensEyePlaneModel}
  \uses{def:physics:phyx-mini-0044:phyxminiproblems-problemphyxmini0044-lengthincentimeters, def:physics:phyx-mini-0044:phyxminiproblems-problemphyxmini0044-contactlenscorrectionsetup}
  The contact lens is modeled at the eye plane, and it must make the desired object appear at the unaided near point. The second equality implements the Cartesian sign convention for an image on the object's side of the lens
\end{definition}
\begin{proof}
  This is the declared model definition of Uses Contact Lens Eye Plane Model.
\end{proof}
\begin{definition}[Has Depicted Sign Configuration]
  \label{def:physics:phyx-mini-0044:phyxminiproblems-problemphyxmini0044-hasdepictedsignconfiguration}
  \lean{PhyXMiniProblems.ProblemPhyXMini0044.HasDepictedSignConfiguration}
  \uses{def:physics:phyx-mini-0044:phyxminiproblems-problemphyxmini0044-lengthincentimeters, def:physics:phyx-mini-0044:phyxminiproblems-problemphyxmini0044-contactlenscorrectionsetup}
  The physical sign branch depicted for correcting near vision in a hyperopic eye
\end{definition}
\begin{proof}
  This is the declared model definition of Has Depicted Sign Configuration.
\end{proof}
\begin{definition}[Satisfies Signed Thin Lens Equation]
  \label{def:physics:phyx-mini-0044:phyxminiproblems-problemphyxmini0044-satisfiessignedthinlensequation}
  \lean{PhyXMiniProblems.ProblemPhyXMini0044.SatisfiesSignedThinLensEquation}
  \uses{def:physics:phyx-mini-0044:phyxminiproblems-problemphyxmini0044-lengthincentimeters, def:physics:phyx-mini-0044:phyxminiproblems-problemphyxmini0044-contactlenscorrectionsetup}
  The signed Gaussian thin-lens equation 1/f = 1/s + 1/s', written in the single centimeter unit system used by the diagram
\end{definition}
\begin{proof}
  This is the declared model definition of Satisfies Signed Thin Lens Equation.
\end{proof}
\begin{definition}[Answer Choice]
  \label{def:physics:phyx-mini-0044:phyxminiproblems-problemphyxmini0044-answerchoice}
  \lean{PhyXMiniProblems.ProblemPhyXMini0044.AnswerChoice}
  The four focal-length readouts printed beside the answer choices, in centimeters
\end{definition}
\begin{proof}
  This is the declared model definition of Answer Choice.
\end{proof}
\begin{definition}[answer Focal Length In Centimeters]
  \label{def:physics:phyx-mini-0044:phyxminiproblems-problemphyxmini0044-answerfocallengthincentimeters}
  \lean{PhyXMiniProblems.ProblemPhyXMini0044.answerFocalLengthInCentimeters}
  \uses{def:physics:phyx-mini-0044:phyxminiproblems-problemphyxmini0044-answerchoice}
  The displayed whole-centimeter focal length for each answer choice
\end{definition}
\begin{proof}
  This is the declared model definition of answer Focal Length In Centimeters.
\end{proof}
\begin{definition}[Matches Answer To Nearest Centimeter]
  \label{def:physics:phyx-mini-0044:phyxminiproblems-problemphyxmini0044-matchesanswertonearestcentimeter}
  \lean{PhyXMiniProblems.ProblemPhyXMini0044.MatchesAnswerToNearestCentimeter}
  \uses{def:physics:phyx-mini-0044:phyxminiproblems-problemphyxmini0044-lengthquantity, def:physics:phyx-mini-0044:phyxminiproblems-problemphyxmini0044-lengthincentimeters, def:physics:phyx-mini-0044:phyxminiproblems-problemphyxmini0044-answerchoice, def:physics:phyx-mini-0044:phyxminiproblems-problemphyxmini0044-answerfocallengthincentimeters}
  Agreement between an exact focal length and a displayed nearest-centimeter answer
\end{definition}
\begin{proof}
  This is the declared model definition of Matches Answer To Nearest Centimeter.
\end{proof}
\begin{lemma}[reciprocal Focal Length eq three hundredths]
  \label{lem:physics:phyx-mini-0044:phyxminiproblems-problemphyxmini0044-reciprocalfocallength-eq-three-hundredths}
  \lean{PhyXMiniProblems.ProblemPhyXMini0044.reciprocalFocalLength_eq_three_hundredths}
  \uses{def:physics:phyx-mini-0044:phyxminiproblems-problemphyxmini0044-lengthincentimeters, def:physics:phyx-mini-0044:phyxminiproblems-problemphyxmini0044-contactlenscorrectionsetup, def:physics:phyx-mini-0044:phyxminiproblems-problemphyxmini0044-matchesprimaryfigure, def:physics:phyx-mini-0044:phyxminiproblems-problemphyxmini0044-satisfiessignedthinlensequation}
  The two signed distances in the primary figure make the reciprocal focal length exactly 3/100 cm⁻¹ by the thin-lens equation
\end{lemma}
\begin{proof}
  The conclusion follows from the governing relations and figure readouts named in the statement of reciprocal Focal Length eq three hundredths.
\end{proof}
% --- Archon declaration topology end ---
% --- Archon physics formalization source end ---
